\documentclass[10pt]{article}
% Shared LaTeX style for MIT 18.06 exam revision notes.
% Compile topic files with Tectonic, XeLaTeX, or LuaLaTeX.

\usepackage[a4paper,margin=0.72in]{geometry}
\usepackage{fontspec}
\usepackage{amsmath,amssymb,mathtools}
\usepackage{booktabs,array,tabularx}
\usepackage{enumitem}
\usepackage{titlesec}
\usepackage{fancyhdr}
\usepackage{microtype}
\usepackage{xcolor}

\setmainfont{Times New Roman}
\setsansfont{Arial}

\setlength{\parindent}{0pt}
\setlength{\parskip}{3.6pt}
\setlength{\abovedisplayskip}{6pt}
\setlength{\belowdisplayskip}{6pt}
\setlength{\abovedisplayshortskip}{4pt}
\setlength{\belowdisplayshortskip}{4pt}
\renewcommand{\arraystretch}{1.08}
\setlength{\tabcolsep}{4.5pt}

\titleformat{\section}
  {\normalfont\large\bfseries}
  {\thesection}
  {0.55em}
  {}
\titlespacing*{\section}{0pt}{10pt}{3pt}

\titleformat{\subsection}
  {\normalfont\normalsize\bfseries}
  {\thesubsection}
  {0.5em}
  {}
\titlespacing*{\subsection}{0pt}{7pt}{2pt}

\setlist[itemize]{leftmargin=1.35em,itemsep=1pt,topsep=2pt,parsep=0pt}
\setlist[enumerate]{leftmargin=1.55em,itemsep=1pt,topsep=2pt,parsep=0pt}

\pagestyle{fancy}
\fancyhf{}
\fancyfoot[C]{\scriptsize\color{black!45}MIT 18.06 Linear Algebra Exam Notes --- \thepage}
\renewcommand{\headrulewidth}{0pt}
\renewcommand{\footrulewidth}{0pt}

\newcommand{\notesubtitle}[1]{%
  {\centering\itshape #1\par}
  \vspace{2pt}
}

\newcommand{\source}[1]{%
  {\centering\small #1\par}
  \vspace{6pt}
}

\newcommand{\labelnote}[2]{%
  \par\smallskip
  \noindent\textbf{#1.}\quad #2
  \par\smallskip
}

\newcommand{\tightlabel}[1]{%
  \par\smallskip
  \noindent\textbf{#1.}\quad
}

\newcolumntype{Y}{>{\raggedright\arraybackslash}X}
\newcolumntype{L}[1]{>{\raggedright\arraybackslash}p{#1}}

\newenvironment{notetable}
  {\small\begin{tabularx}{\linewidth}}
  {\end{tabularx}}


\begin{document}

\begin{center}
{\LARGE Permutations, Subspaces, Rank, and Dot Products\par}
\end{center}
\notesubtitle{Concise revision notes from the 30 Dec handwritten matrix-operations review.}
\source{Based on handwritten MIT 18.06 revision notes.}

\section{Core Idea}

\labelnote{Core idea}{Early linear algebra questions often reduce to three checks: what operations preserve structure, whether a set is a subspace, and how rank controls solution behavior.}

This note collects matrix-operation facts used repeatedly later: permutation matrices, subspaces, rank cases, row/column relationships, dot products, and skew-symmetry.

\section{Permutation Matrices}

A permutation matrix reorders rows or columns. It is orthogonal:
\[
P^{T}P=I,
\qquad
P^{-1}=P^{T}.
\]
Multiplying on the left reorders rows:
\[
PA.
\]
Multiplying on the right reorders columns:
\[
AP.
\]

\labelnote{Exam memory}{A permutation matrix looks like the identity matrix with its rows or columns rearranged.}

\section{Subspace Test}

A set \(S\) is a subspace when:
\begin{enumerate}
  \item \(0\in S\);
  \item if \(u,v\in S\), then \(u+v\in S\);
  \item if \(u\in S\), then \(cu\in S\).
\end{enumerate}

The nullspace is a subspace because if
\[
Au=0,
\qquad
Av=0,
\]
then
\[
A(u+v)=0,
\qquad
A(cu)=0.
\]

\labelnote{Common mistake}{A set defined by \(Ax=b\) with \(b\ne0\) is usually not a subspace, because it does not contain the zero vector.}

\section{Rank Cases}

For \(A\in\mathbb{R}^{m\times n}\), rank \(r\) controls solution behavior.

\begin{center}
\small
\begin{tabularx}{\linewidth}{@{}L{0.20\linewidth}Y Y@{}}
\toprule
Rank case & Meaning & Exam consequence \\
\midrule
\(r=n\) & Full column rank. & At most one solution to \(Ax=b\). \\
\(r=m\) & Full row rank. & \(Ax=b\) is solvable for every \(b\in\mathbb{R}^m\). \\
\(r=m=n\) & Full rank square. & Exactly one solution for every \(b\). \\
\(r<m\) and \(r<n\) & Not full rank. & Zero or infinitely many solutions. \\
\bottomrule
\end{tabularx}
\end{center}

\section{Row Space and Column Space}

Row space and column space have the same dimension:
\[
\dim C(A)=\dim C(A^{T})=\operatorname{rank}(A).
\]
But they usually live in different ambient spaces:
\[
C(A)\subseteq\mathbb{R}^{m},
\qquad
C(A^{T})\subseteq\mathbb{R}^{n}.
\]

\labelnote{Exam memory}{Same dimension does not mean the same space. Check where the vectors live.}

\section{Dot Product}

The dot product is
\[
x\cdot y=x^{T}y.
\]
It measures alignment:
\[
x^{T}y=0
\]
means the vectors are perpendicular.

It also interacts cleanly with matrix multiplication:
\[
(Px)\cdot(Py)=x\cdot y
\]
for a permutation or orthogonal matrix \(P\).

\section{Skew-Symmetric Matrices}

A matrix is skew-symmetric when
\[
A^{T}=-A.
\]
The diagonal entries must be zero, because
\[
a_{ii}=-a_{ii}
\quad\Rightarrow\quad
a_{ii}=0.
\]
Off-diagonal entries come in opposite pairs.

\section{Common Mistakes}

\begin{center}
\small
\begin{tabularx}{\linewidth}{@{}L{0.40\linewidth}Y@{}}
\toprule
Mistake & Correct exam habit \\
\midrule
Treating \(Ax=b\) as a subspace when \(b\ne0\). & Check whether the zero vector is included. \\
Confusing row operations with column operations. & Left multiply changes rows; right multiply changes columns. \\
Assuming equal dimensions imply equal spaces. & Check the ambient spaces. \\
Forgetting the full-rank square case. & \(r=m=n\) gives existence and uniqueness. \\
Putting nonzero entries on a skew-symmetric diagonal. & The diagonal must be zero. \\
\bottomrule
\end{tabularx}
\end{center}

\section{Final Exam Checklist}

\tightlabel{Checklist}
\begin{enumerate}
  \item Identify whether multiplication is on the left or right.
  \item Check \(P^{T}P=I\) for permutation matrices.
  \item Test subspaces by zero vector, addition, and scalar multiplication.
  \item Write the size \(m\times n\) before using rank facts.
  \item Find rank \(r\).
  \item Use full column rank for uniqueness.
  \item Use full row rank for existence.
  \item Use \(x^{T}y=0\) to check perpendicularity.
  \item Check skew-symmetry with \(A^{T}=-A\).
\end{enumerate}

\end{document}
